\section{A rational abelian-sector target for arbitrary smooth
projective
cohomology}\label{a-rational-abelian-sector-target-for-arbitrary-smooth-projective-cohomology}

The remaining cross-cycle question in cote 15 can be put against one
\textbf{rational} target determined by the Frobenius polynomials of the
specified variety. The rational orbit and Schur-index obstructions in
CC-AG08 automatically vanish for the actual cohomology of a smooth
projective variety. The reason is crystalline Frobenius at the finite
places and the hard-Lefschetz pairing at the real places. This is a
deduction from those established results and the previously scoped
abelian-sector classification, with no historical-priority claim.

The arithmetic divisibility statement below is unconditional. Its
interpretation inside the rational abelian--Artin--Tate category uses
exactly the all-abelian inputs already recorded in CC-TATE01--02,
AQ-ABELIAN-RECOGNITION, and CC-AG08. No Tate assertion, numerical
realization, or Frobenius semisimplicity is assumed for the specified
variety.

\subsection{Primary inputs and
conventions}\label{primary-inputs-and-conventions}

Let \(k=\mathbf F_q\), \(q=p^a\), and let \(X/k\) be a nonempty smooth
projective scheme, pure of dimension \(d\). Use geometric Frobenius
\(F\) on étale cohomology. For \(0\leq i\leq 2d\), put \[
 P_i(t)=\det(t-F\mid H^i_{\mathrm{et}}
                  (X_{\bar k},\mathbf Q_\ell))\in\mathbf Z[t].
                                                        \tag{CC-XC.R11}
\] This is the monic characteristic polynomial, rather than the
\(\det(1-tF)\) convention. It is independent of \(\ell\ne p\); its roots
are nonzero integral \(q\)-Weil numbers of weight \(i\).

The precise primary results used are:

\begin{itemize}
\tightlist
\item
  Deligne (1974), Theorem 1.6, printed p.~276, gives the finite-field
  Weil theorem in the smooth projective case. Katz and Messing (1974),
  Theorem 1, printed p.~74, identify the characteristic polynomials in a
  reasonable Weil cohomology with the étale ones. Their immediately
  preceding paragraph explicitly includes crystalline cohomology. Their
  Corollary 1(2), p.~75, also transfers hard Lefschetz to those
  theories.
\item
  Put \(K_0=W(k)[1/p]\), and let \(\sigma\) be Witt Frobenius, of order
  \(a\). The rational crystalline cohomology
  \(C_i=H^i_{\mathrm{cris}}(X/W(k))[1/p]\) is a finite-dimensional
  \(K_0\)-space with an invertible \(\sigma\)-semilinear Frobenius
  \(\varphi\). Thus \(\varphi^a\) is \(K_0\)-linear. The Stacks Project,
  Remark 60.24.13, tag 07MY, gives perfectness, and Theorem 60.26.4, tag
  07N5, gives Frobenius invertibility after inverting \(p\). Their
  hypotheses hold for the structure crystal over \((W(k),(p),\gamma)\);
  Situation 60.26.1 and Definition 60.26.2, tags 07N2--07N3, specify the
  Frobenius lift and the nondegenerate structure.
\item
  Deligne (1980), Theorem 4.1.1, printed p.~217, gives hard Lefschetz
  for a smooth projective variety of pure dimension over an
  algebraically closed field and an ample line bundle. We use its
  ordinary Poincaré-duality consequence, with an ample line bundle
  already defined over \(k\).
\item
  The rational simple-object index formula is the previously scoped
  CC-AG08: Milne (1994), corrected author version of 2 November 2019,
  Proposition 2.8, p.~14, and Theorems 2.15--2.16, p.~18. Here
  \(\operatorname{End}\) means the endomorphism algebra of the
  cohomological motive, with the geometric-Frobenius convention
  specified on Milne's printed p.~17.
\end{itemize}

For clarity, Katz--Messing state their argument for an absolutely
irreducible smooth projective variety. The polynomial equality needed
here also follows for our \(X\). Partition its geometric connected
components into Frobenius orbits. For an orbit of length \(s\), one
component is defined over \(\mathbf F_{q^s}\) and is geometrically
connected and smooth, hence geometrically irreducible. In either
cohomology theory, after extending coefficients, the orbit contributes
the direct sum of the \(s\) component spaces, cyclically permuted by
\(F\), with \(F^s\) on a chosen component equal to its
\(\mathbf F_{q^s}\)-Frobenius. In the crystalline statement this means
the linear base change of \(\varphi^a\), whose \(s\)-th power agrees
with the \(\varphi^{as}\) for the enlarged constant field; it does not
mean a semilinear extension acting on the new coefficient field. A
cyclic block operator \(B\) satisfies \[
 \det(t-B)=\det(t^s-B^s\mid\text{one component});
                                                        \tag{CC-XC.R12}
\] one sees this by using \(B\) to identify successive component spaces,
which gives the usual block companion matrix. Katz--Messing on the
chosen component proves the equality on the orbit, and then on their
direct sum. The crystalline coefficient extension used here is
proper-smooth base change: Stacks, Remark 60.24.9, tag 07MU, modulo
\(p^n\), followed by perfect \(p\)-adic completion as in 07MY. We
therefore have, without assuming semisimplicity, \[
 \det(t-\varphi^a\mid C_i)=P_i(t).
                                                        \tag{CC-XC.R13}
\]

\subsection{CC-XC11: automatic Schur-index divisibility and its rational
target}\label{cc-xc11-automatic-schur-index-divisibility-and-its-rational-target}

Factor the monic rational polynomial as \[
 P_i(t)=\prod_f f(t)^{m_f},
                                                        \tag{CC-XC.R14}
\] where the \(f\) are distinct monic irreducible polynomials in
\(\mathbf Q[t]\) and \(m_f>0\). Choose a root \(\alpha_f\), put
\(K_f=\mathbf Q(\alpha_f)\), and normalize every \(p\)-adic valuation by
\(v_p(p)=1\). Define \[
 d_f=\operatorname{lcm}_v
 \operatorname{den}\begin{cases}
 i/2,&v\text{ real},\\[2pt]
 -[K_{f,v}:\mathbf Q_p]\,v_p(\alpha_f)/a,&v\mid p,\\[2pt]
 0,&\text{otherwise}.
 \end{cases}
                                                        \tag{CC-XC.R15}
\] The entries are taken in \(\mathbf Q/\mathbf Z\), and the denominator
of zero is one. These are exactly the local invariants and Schur index
of the rational abelian-sector simple object labelled by the Galois
orbit of \(\alpha_f\).

\textbf{Theorem.} For every factor in (CC-XC.R14), \[
                         d_f\mid m_f.                  \tag{CC-XC.R16}
\] Consequently there is a unique rational isomorphism class
\(M_i(X)\in\mathcal A_q\), pure of weight \(i\), with polynomial
\(P_i(t)\). If \(S_f\) is the rational simple with label \(f\), it is \[
 M_i(X)=\bigoplus_f S_f^{\oplus m_f/d_f},\qquad
 M(X)=\bigoplus_{i=0}^{2d}M_i(X).
                                                        \tag{CC-XC.R17}
\] For every \(\ell\ne p\), its realization is isomorphic to the
semisimplification of the graded Frobenius module \(H^*(X_{\bar k},
\mathbf Q_\ell)\). The target admits a rational presentation \[
 M(X)=(T,e),\qquad
 T=\bigoplus_{\nu=1}^{s}h(Y_\nu),
                                                        \tag{CC-XC.R18}
\] where every \(Y_\nu/k\) is an abelian variety or a finite étale
scheme and \(e\) is a rational numerical projector. By the established
numerical--homological equality on this sector, it is a projector in
every corresponding homological realization. No Tate twist is needed in
(CC-XC.R18).

The unique isomorphism class in (CC-XC.R17) is a spectral target. The
statement does not identify \(h(X)\) with it.

\subsubsection{Complete proof}\label{complete-proof}

\textbf{The \(p\)-adic denominators.} Fix \(f\), and fix a place
\(v\mid p\) of \(K_f\). The place corresponds to an irreducible factor
\(f_v\in\mathbf Q_p[t]\) of \(f\), of degree \[
 n_v=\deg f_v=[K_{f,v}:\mathbf Q_p].
\] Set \(F_{\mathrm{cris}}=\varphi^a\) and choose an integer \(B\) at
least \(\dim_{K_0} C_i\). On \(C_i\) take the entire
\(\mathbf Q_p\)-primary block \[
 U_v=\ker\bigl(f_v(F_{\mathrm{cris}})^B\bigr).
                                                        \tag{CC-XC.R19}
\] Its dimension is \(m_f n_v\): after an algebraic closure of the
coefficient field it is precisely the sum of the generalized eigenspaces
for the \(n_v\) distinct roots of \(f_v\), each of algebraic
multiplicity \(m_f\).

It is essential here to take the whole \(f_v\)-block, not an individual
factor over \(K_0\). Since the coefficients of \(f_v\) lie in
\(\mathbf Q_p\), they are fixed by \(\sigma\). Also
\(\varphi F_{\mathrm{cris}}=F_{\mathrm{cris}}\varphi\). Consequently
(CC-XC.R19) is preserved by \(\varphi\) and by its inverse. In a
\(K_0\)-basis of \(U_v\), write \(\varphi=A\sigma\) with
\(A\in\mathrm{GL}_{m_fn_v}(K_0)\). Then \[
 \det(F_{\mathrm{cris}}\mid U_v)
   =\det A\cdot\sigma(\det A)\cdots\sigma^{a-1}(\det A)
   =N_{K_0/\mathbf Q_p}(\det A).
                                                        \tag{CC-XC.R20}
\] The extension \(K_0/\mathbf Q_p\) is unramified, so \(v_p(\det A)\)
is an integer and \(\sigma\) preserves the valuation. All roots of the
irreducible \(f_v\) have the same \(p\)-adic valuation as \(\alpha_f\)
at \(v\). Taking valuations of (CC-XC.R20) gives \[
 m_f n_v\,v_p(\alpha_f)=a\,v_p(\det A),\qquad
 m_f\,\frac{[K_{f,v}:\mathbf Q_p]v_p(\alpha_f)}a
        \in\mathbf Z.                                 \tag{CC-XC.R21}
\] Thus each finite-place denominator in (CC-XC.R15) divides \(m_f\).
This argument used the characteristic polynomial, with its full
algebraic multiplicities, and never discarded a possible Jordan block.

\textbf{A linear algebra fact for the real places.} Suppose \(V\) is a
finite-dimensional vector space over an algebraically closed
characteristic-zero field, \(B\) a nondegenerate alternating form, and
\(F\) an invertible operator satisfying \[
 B(Fx,Fy)=c\,B(x,y),\qquad c\ne0.
                                                        \tag{CC-XC.R22}
\] If \(V_\lambda\) denotes the generalized \(\lambda\)-eigenspace, then
\[
 B(V_\lambda,V_\mu)=0\quad\hbox{unless }\lambda\mu=c.
                                                        \tag{CC-XC.R23}
\] Indeed the adjoint of \(F\) is \(cF^{-1}\). If \(\lambda\mu\ne c\),
then \(cF^{-1}-\lambda\) is invertible on \(V_\mu\). For
\(x\in V_\lambda\), choose \(N\) with \((F-\lambda)^N x=0\). Every
\(y\in V_\mu\) equals \((cF^{-1}-\lambda)^N z\) for some \(z\in V_\mu\),
and the adjoint identity gives \[
 B(x,y)=B((F-\lambda)^N x,z)=0.
 \] The generalized eigenspaces sum to all of \(V\), so the
nondegeneracy of \(B\) implies that its pairing between \(V_\lambda\)
and \(V_{c/\lambda}\) is perfect. In particular, if \(\lambda^2=c\), the
restriction to \(V_\lambda\) is a nondegenerate alternating form, and \[
                         \dim V_\lambda\text{ is even}. \tag{CC-XC.R24}
\] This proof does not require \(F\) to be diagonalizable.

\textbf{Odd-degree cohomology.} First suppose \(i\le d\) is odd. Choose
an ample line bundle on \(X\) over \(k\), and view its first Chern class
as \(L\in H^2(X_{\bar k},\mathbf Q_\ell)\) after trivializing the Tate
twist. Then \(F(L)=qL\). Hard Lefschetz and Poincaré duality show that
\[
 B_i(x,y)=\operatorname{Tr}(x\smile y\smile L^{d-i})
                                                        \tag{CC-XC.R25}
\] is nondegenerate and alternating on \(H^i\). The trace means the sum
of the component traces when \(X_{\bar k}\) is disconnected; this
pairing is still perfect. Frobenius acts on a component's top cohomology
by \(q^d\) and permutes the components. Therefore \[
 \begin{aligned}
 q^{d-i}B_i(Fx,Fy)
 &=\operatorname{Tr}
       \bigl(F(x\smile y\smile L^{d-i})\bigr)\\
 &=q^d B_i(x,y),
 \end{aligned}
 \qquad B_i(Fx,Fy)=q^iB_i(x,y).                         \tag{CC-XC.R26}
\] Apply (CC-XC.R24) after extending coefficients to
\(\overline{\mathbf Q}_\ell\). Every occurrence of \(\lambda=q^{i/2}\)
or \(\lambda=-q^{i/2}\) has even algebraic multiplicity.

For odd \(i>d\), put \(j=2d-i<d\). Hard Lefschetz gives the isomorphism
\[
 L^{i-d}:H^j\xrightarrow{\sim}H^i,\qquad
 F_i L^{i-d}=q^{i-d}L^{i-d}F_j.                         \tag{CC-XC.R27}
\] It carries generalized eigenspaces for \(\beta\) to those for
\(q^{i-d}\beta\), preserving their dimensions. The real eigenvalues
\(\pm q^{i/2}\) therefore correspond to \(\pm q^{j/2}\), whose
multiplicities are even by the preceding case.

Now if \(K_f\) has a real embedding, a real conjugate of \(\alpha_f\)
must be one of \(\pm q^{i/2}\), by its Weil absolute value. Its
irreducible polynomial divides \(t^2-q^i\). When this splits over
\(\mathbf Q\), each relevant linear factor has even exponent in odd
degree. When it is irreducible, the centre is the real quadratic field
\(\mathbf Q(q^{i/2})\); each of its two roots has generalized eigenspace
dimension \(m_f\), so \(m_f\) is again even. For even \(i\), the real
invariant \(i/2\) is zero and imposes no condition. Thus every
real-place denominator in (CC-XC.R15) also divides \(m_f\).

The other invariants vanish. Taking their least common multiple proves
(CC-XC.R16).

\textbf{The rational motive and its effective presentation.} Each root
\(\alpha\) is a \(p\)-unit: both \(\alpha\) and its complex conjugate
\(q^i/\alpha\) are algebraic integers, so their valuations at a prime
away from \(p\) are nonnegative and sum to zero. Thus the roots belong
to the Weil-label group of CC-AG08. They have constant multiplicity
\(m_f\) on each rational Galois orbit. Equation (CC-XC.R16) is exactly
the second divisibility condition in CC-AG08. That theorem gives
(CC-XC.R17), including uniqueness of the rational isomorphism class. Its
realization is semisimple, with characteristic polynomial \(P_i\). Over
any characteristic-zero field, two semisimple operators with the same
characteristic polynomial are isomorphic as modules over the polynomial
ring: decompose them by irreducible factors, whose simple modules are
the corresponding residue fields. This proves the assertion about every
\(\ell\)-adic semisimplification.

For the presentation without twists, consider first a root \(\alpha\) in
positive degree \(i\). Choose \(\gamma\) with \(\gamma^i=\alpha\). It is
integral, and every complex conjugate has absolute value \(q^{1/2}\), so
it is a weight-one integral \(q\)-Weil number. Honda's Main Theorem
supplies an abelian variety \(B/k\) with \(\gamma\) on \(H^1\). The
Künneth summand \(H^1(B)^{\otimes i}\) of \(H^i(B^i)\) then has
eigenvalue \(\alpha\). In degree zero, every root is a root of unity,
and a root of order \(n\) occurs in
\(h(\operatorname{Spec}\mathbf F_{q^n})\).

Thus each simple \(S_f\) appearing in (CC-XC.R17) occurs in the rational
semisimple decomposition of the cohomology motive of one such abelian
variety or finite étale scheme. Indeed its realization contains a label
in the orbit \(f\), and the classification CC-AG08 identifies the unique
rational simple with that orbit. Take enough copies to accommodate the
desired integer number \(m_f/d_f\) of each rational simple.
Semisimplicity then supplies a rational summand projector on their
finite direct sum \(T\). On this sector numerical and homological
equivalence agree, so a rational cycle representative of the projector
gives (CC-XC.R18) in every realization. This completes the proof.
\(\square\)

The index \(d_f\) and the abstract rational isomorphism class (CC-XC.R17)
can be computed from the factorizations of \(P_i\), number-field local
degrees and valuations. This assertion does not say that those
polynomials alone print equations for the varieties \(Y_\nu\) or cycles
representing \(e\); a presentation in the already understood abelian
sector is additional concrete data to be chosen. Once chosen, it remains
fixed throughout the cross-cycle search below.

\subsection{CC-XC12: one fixed rational target and a finite cycle
certificate}\label{cc-xc12-one-fixed-rational-target-and-a-finite-cycle-certificate}

Fix a rational presentation (CC-XC.R18) of the target just proved to
exist. Write \(r_\nu=\dim Y_\nu\), treating components of a finite étale
\(Y_\nu\) as dimension zero. For rational cycles \[
 z_\nu\in\mathrm{CH}^d(X\times Y_\nu)_{\mathbf Q},\qquad
 w_\nu\in\mathrm{CH}^{r_\nu}(Y_\nu\times X)_{\mathbf Q},
                                                        \tag{CC-XC.R28}
\] let \[
 f=e\circ z:h(X)\longrightarrow M(X),\qquad
 g=w\circ e:M(X)\longrightarrow h(X).
                                                        \tag{CC-XC.R29}
\] These expressions refer first to correspondences; the numerical
projector \(e\) is used through its homological class on the abelian
sector. The codimensions in (CC-XC.R28) give degree-zero cohomological
maps in both directions.

\textbf{Theorem.} The following conditions are equivalent.

\begin{enumerate}
\def\labelenumi{\arabic{enumi}.}
\tightlist
\item
  For some \(\ell\ne p\), the object \(h_{\ell,E}(X)\), with
  \(E=\overline{\mathbf Q}_\ell\), belongs to the abelian sector of the
  homological category used in CC-XC02.
\item
  For some such \(\ell\), the determinant is nonzero on the image of the
  single projected cross-cycle map \[
  \bigoplus_\nu
      \mathrm{CH}^d(X\times Y_\nu)_{\mathbf Q}\otimes E
    \longrightarrow
      \operatorname{Hom}_{E[F]}
      \bigl(H^*(X_{\bar k},E),R_E(M(X))\bigr),\quad
        z\longmapsto R_E(ez).
                                                       \tag{CC-XC.R30}
  \]
\item
  There are rational cycles (CC-XC.R28) for which \[
      \det(gf\mid H^i(X_{\bar k},\mathbf Q_\ell))\ne0
         \quad\text{for every }0\le i\le2d.
                                                       \tag{CC-XC.R31}
  \] Here the determinants are rational numbers independent of
  \(\ell\ne p\), so the condition can be checked at no chosen
  cohomological prime.
\item
  Rational correspondences give mutually inverse isomorphisms between
  \(h(X)\) and this fixed \(M(X)\) in every \(\ell\)-adic homological
  realization, \(\ell\ne p\).
\end{enumerate}

The certificates in (3) are finite and their verification is the
intersection-number procedure of CC-XC05. Enumerating only the cycles
(CC-XC.R28), with \(T\) and \(e\) fixed, is a one-sided procedure that
halts exactly when these conditions hold. It does not decide
nonmembership by failing to halt.

\subsubsection{Complete proof}\label{complete-proof-1}

CC-XC11 shows that the fixed \(R_E(M(X))\) is precisely the semisimple
target required in CC-XC02, with the same dimension in each degree. Its
proof uses only a presentation by abelian varieties and finite étale
schemes, which (CC-XC.R18) provides. CC-XC02 therefore applies verbatim
to (CC-XC.R30): (1) and (2) are equivalent. If they hold, an invertible
projected map has an algebraic inverse in the \(E\)-coefficient
homological category, by the duality argument proved there.

Write such an inverse pair as finite sums of rational cycle
representatives with coefficients in \(E\): \[
 f_E=\sum_{u=1}^{b}c_u\,e z_u,\qquad
 g_E=\sum_{v=1}^{c}d_v\,w_v e.
 \] The projector is already rational. The product \[
 D(x,y)=\prod_{i=0}^{2d}
       \det\!\left(
       \left.\sum_{u,v}x_u y_v\,w_v e z_u
       \right|H^i(X_{\bar k},E)\right)
                                                        \tag{CC-XC.R32}
 \] is a polynomial in finitely many variables over \(E\), nonzero at
the coefficients of \(f_E,g_E\). A nonzero polynomial over an extension
field of \(\mathbf Q\) cannot vanish on every rational tuple: induct on
the number of variables, using the infinitely many rational values for
the last variable. Choose such a tuple. It gives rational \(f,g\)
satisfying (CC-XC.R31) at the selected prime.

For these rational correspondences, CC-XC05 supplies the
prime-independent rational characteristic polynomial of \(gf\) in every
degree. Explicitly, the rational Frobenius polynomials \(P_i\) are
pairwise coprime by their weights; the Chinese remainder theorem, using
the full characteristic polynomials \(P_i\) including their repeated
factors, gives rational polynomials in Frobenius acting as cohomological
projectors. Their construction needs no semisimplicity. The intersection
trace of the composition of such a projector with \((gf)^n\), including
its parity sign, gives the trace in the single degree \(i\). Newton's
identities then give that degree's characteristic polynomial. All these
intersection degrees are rational and independent of \(\ell\), as proved
in CC-XC05. Therefore the nonzero determinants just obtained are nonzero
in every realization. This proves (3).

Conversely suppose (3). On every degree, Cayley--Hamilton expresses the
inverse of \(gf\) as a polynomial in \(gf\) with rational coefficients:
the constant term of its characteristic polynomial is nonzero. Combine
these inverse polynomials using the rational cohomological projectors
just described. This gives a rational correspondence \(u\) acting as
\((gf)^{-1}\) in every degree and every \(\ell\ne p\). Set \(g'=ug\).
Then \(g'f=1\) on every realization of \(X\). Since the dimensions of
\(X\) and \(M(X)\) agree in each degree by CC-XC11, \(f\) is an
isomorphism there and \(fg'=1\) on \(M(X)\) as well. Identities of these
linear actions are identities in each homological category, by its
definition. Thus (4) holds. Condition (4) implies (1), completing the
equivalences.

Finally, rational algebraic cycles over the fixed finite field, with
rational coefficients, can be enumerated as in CC-XC07. The fixed
numerical projector is represented by a fixed rational matrix of cycles
on the products \(Y_\nu\times Y_\mu\). For each pair of lists
(CC-XC.R28), the finite compositions, intersections, traces, and Newton
identities are exactly the certificate verifier CC-XC05. A successful
certificate proves (3). If (1) holds, the finite rational specialization
above provides some certificate in the enumeration, so the procedure
eventually finds it. No step asserts termination when (1) fails.
\(\square\)

\subsection{What this removes and what
remains}\label{what-this-removes-and-what-remains}

For arbitrary smooth projective \(X/\mathbf F_q\), neither Galois-orbit
multiplicities nor local Schur indices are an additional rational-target
obstruction: (CC-XC.R16) proves their required divisibilities. The
spectral target is uniquely fixed before searching for
cross-correspondences. The remaining condition is a finite positive
certificate on \(X\) and that target, rather than a search for a target
isomorphism class.

The proof has not made Frobenius on \(H^*(X)\) semisimple:
semisimplicity follows only after the determinant certificate, because
then \(H^*(X)\) is isomorphic to the semisimple abelian-sector
realization. Nor has it supplied the cross-cycles for all \(X\). The
arbitrary-variety Tate/effectivity question therefore retains this
explicit cycle-existence gate. Projectivity is used in the real-place
argument through hard Lefschetz; the proof is not asserted for arbitrary
smooth nonproper schemes or without that pairing. The already proved
nonexistence of a \(\mathbf Q\)-valued fibre functor on the whole
rational abelian category also remains unchanged.

\subsection{APA bibliography and exact
locators}\label{apa-bibliography-and-exact-locators}

Deligne, P. (1974). La conjecture de Weil: I. \emph{Publications
Mathématiques de l'IHÉS, 43}, 273--307. Theorem 1.6, p.~276.
https://www.numdam.org/item/PMIHES\_1974\_\_43\_\_273\_0.pdf

Deligne, P. (1980). La conjecture de Weil: II. \emph{Publications
Mathématiques de l'IHÉS, 52}, 137--252. Theorem 4.1.1, p.~217; the
paired odd-Betti-number consequence is Corollary 4.1.5, p.~218.
https://publications.ias.edu/sites/default/files/Number40.pdf

Honda, T. (1968). Isogeny classes of abelian varieties over finite
fields. \emph{Journal of the Mathematical Society of Japan, 20}(1--2),
83--95. Main Theorem, p.~83. https://doi.org/10.2969/jmsj/02010083

Katz, N. M., \& Messing, W. (1974). Some consequences of the Riemann
hypothesis for varieties over finite fields. \emph{Inventiones
Mathematicae, 23}, 73--77. Theorem 1, p.~74; Corollary 1, p.~75; Theorem
2, p.~76, for the cohomological-projector input already used in CC-XC05.
https://web.math.princeton.edu/\textasciitilde nmk/old/katzmessing.pdf

Milne, J. S. (1994). Motives over finite fields. In U. Jannsen, S. L.
Kleiman, \& J.-P. Serre (Eds.), \emph{Motives} (Proceedings of Symposia
in Pure Mathematics, Vol. 55, Part 1, pp.~401--459). American
Mathematical Society. Corrected author version, 2 November 2019:
Propositions 2.4, 2.6, 2.8, pp.~13--14; endomorphism convention, p.~17;
Theorems 2.15--2.16, p.~18. The full-Tate hypothesis in its Section 2 is
used here only through the already justified abelian-sector restriction
CC-AG08. https://www.jmilne.org/math/xnotes/Seattle1.pdf

The Stacks Project Authors. (2026). \emph{The Stacks Project}. Remark
60.24.9 (base change), tag 07MU; Remark 60.24.13 (perfectness), tag
07MY; Situation 60.26.1, tag 07N2; Definition 60.26.2, tag 07N3; Theorem
60.26.4 (Frobenius after inverting \(p\)), tag 07N5.
https://stacks.math.columbia.edu/tag/07MU ;
https://stacks.math.columbia.edu/tag/07MY ;
https://stacks.math.columbia.edu/tag/07N2 ;
https://stacks.math.columbia.edu/tag/07N3 ;
https://stacks.math.columbia.edu/tag/07N5

Proved internal inputs: CC-AG08 (rational abelian-sector
classification), CC-XC02 (one projected determinant criterion), and
CC-XC05--07 (finite rational-cycle certificates and their one-sided
enumeration). The external computability and intersection-theory inputs
to the last three results are retained with their exact locators in
research/consequences\_criteria\_cycle\_certificates.tex.
